% Primary-source checks, with their exact version and coefficient scope.
\section{Two Stieltjes index and coefficient corrections}
\label{dgp:audit:section}
These checks concern the identified displays in the inspected source PDFs.
They also illustrate why the coefficient convention and the initial
summation index must be fixed before comparing harmonic spectral identities.

\subsection{The initial term in the Hurwitz--Stieltjes limit}
\label{dgp:audit:Can}
On printed page 1 of \emph{On certain harmonic zeta functions},
arXiv:2403.07123v1, dated 11 March 2024 \cite{dgp:CanV1}, the Hurwitz
series begins at $n=0$, but the following unnumbered Stieltjes-limit
display begins at $n=1$. For real $a>0$ and integer $m\ge0$, the correct
unit-coordinate limit is
\begin{equation}\label{dgp:audit:correct-limit}
 \gamma_m(a)=\lim_{N\to\infty}\left[
   \sum_{n=0}^N\frac{\log^m(n+a)}{n+a}
        -\frac{\log^{m+1}(N+a)}{m+1}\right].
\end{equation}
Equivalently, retain the lower index one and add $\log^m(a)/a$.
\begin{proof}
The standard sum--integral subtraction for the Laurent expansion of
$\zeta(s,a)$ gives \eqref{dgp:audit:correct-limit}; changing the upper
endpoint by one contributes $O((1+\log^m N)/N)$ and does not change
the limit. The recurrence
$\zeta(1+u,a+1)=\zeta(1+u,a)-a^{-1-u}$ gives, coefficient by coefficient,
\[
 \gamma_m(a+1)=\gamma_m(a)-\frac{\log^m a}{a}.
\]
The sum starting at one is the sum for $a+1$, hence its printed limit is
$\gamma_m(a+1)$, not $\gamma_m(a)$.
\end{proof}
At $m=0,a=1$, the exact counterexample is
\[
 \lim_{N\to\infty}\{H_{N+1}-1-\log(N+1)\}=\gamma-1.
\]
At $a=1,m>0$ the omitted term vanishes; testing only those cases would
miss the index defect. This correction is specifically to the inspected
version-one introductory display.

\subsection{The sign of the height-one linear coefficient}
\label{dgp:audit:Young}
Young's published \emph{Series of height one multiple zeta functions},
\emph{Integers} 24 (2024), A43, printed page 12 \cite{dgp:Young}, describes
the symbol in Corollary 3 as the linear coefficient. Its displayed
integral expression instead gives the negative of the literal coefficient
$L_d=[u]Z(1+u,1^{d-1})$ in our convention. The ordinary Stieltjes value in
Example 5 is consistent with the conventional $\gamma_1=-L_1$.

The sign can be checked at every depth without a decimal fit. For
$0\le r\le1$ and $s>1$, define the classical Barnes series
\[
 \mathcal B_r(s)=\sum_{n\ge0}\frac{(r)_n}{n!}(n+1)^{-s}.
\]
At $r=0$ only the initial term remains. Put
\[
 U_d=\int_0^1(1-r)\mathcal B_r(d+2)\,\dd r,
 \qquad V_d=\int_0^1\mathcal B_r(d+1)\,\dd r.
\]
The right side of the inspected Corollary 3 is $U_d-V_d-\alpha_{d+1}$.
Using $(r)_{n+1}=(r+n)(r)_n$ and termwise integration gives
\[
 U_d-V_d=-\sum_{n\ge1}\frac{\rho_n}{n^{d+2}},
\]
where $\rho_n$ is defined in Proposition~\ref{dgp:powers:prop:Cauchy}.
All relevant series converge absolutely because $(r)_n/n!\le1$ on
$0\le r\le1$ and $d\ge1$. The proved Cauchy coefficient formula therefore
shows exactly that
\begin{equation}\label{dgp:audit:Young-sign}
 U_d-V_d-\alpha_{d+1}=-L_d.
\end{equation}
Thus one may keep the displayed expression and describe the negative
linear coefficient, or negate it when naming the literal $[u]$ coefficient.
At depth one it gives $\gamma_1$, whereas the literal linear coefficient
of $\zeta(1+u)$ is $-\gamma_1$. The depth-one discrepancy is nonzero rigorously: the exact rational
certificate in \path{verification/digamma-coordinate-source-audit.json}
gives
\[
 -\frac3{40}<\gamma_1<-\frac7{100}<0.
\]
Here is its analytic error budget. For $f(x)=\log(x)/x$ and $N=10$,
one Euler--Maclaurin correction gives
\[
 \gamma_1=\sum_{n=1}^{N}\frac{\log n}{n}
   -\frac{\log^2N}{2}-\frac{\log N}{2N}
   +\frac{\log N-1}{12N^2}+E_N,
 \qquad |E_N|\le\frac{\log N-1}{12N^2}.
\]
The bound follows from $|\widetilde B_2|/2!\le1/12$ and
$f''(x)=(2\log x-3)/x^3\ge0$ for $x\ge10$; thus
$\int_N^\infty|f''|=-f'(N)$. Each logarithm is enclosed rationally by
128 terms of its atanh series, with the geometric positive-tail bound.
The resulting endpoint comparisons are exact fractions, not floating
sign tests. The separately replayed 100-digit Cauchy/Barnes diagnostic
also confirms the sign; \eqref{dgp:audit:Young-sign} proves the
all-depth coefficient identity. The correction concerns that coefficient description,
with Example 5's ordinary Stieltjes formula retained.
